\section{Checkpoint Selection}
\label{app:checkpoint-selection}

\begin{table}[h]
  \centering
  \scriptsize
  \resizebox{\columnwidth}{!}{%
  \begin{tabular}{ll}
    \toprule
    Method & Candidate values per concept \\
    \midrule
    EMBER (Rome) & $\delta\in\{1,2,5,10,20,50,100,200,500,1000,2000,5000\}$ \\
    EMBER (Baseball, AI) & $\delta\in\{1,2,5,10,20,50,100,200,500\}$ \\
    RMU, RMU+EMBER & layer $\in\{5,6\}$; steering $\in\{\text{mid},\text{high}\}$; $\alpha\in\{10,100\}$ \\
    SNMF & side $\in\{\text{input},\text{output},\text{both}\}$; Rome also judged-both \\
    SNMF+EMBER & ratio selection; side $\in\{\text{input},\text{output},\text{both}\}$ \\
    \bottomrule
  \end{tabular}%
  }
  \caption{Complete candidate grid. Cartesian products are taken within each
  row. EMBER's $\delta$ is edit strength; RMU's $\alpha$ weights retention.}
  \label{tab:hyperparameter-grid}
\end{table}

\subsection{Candidate Grid and Selection Procedure}

We selected checkpoints separately for each combination of concept and
erasure method. The candidate grid contained 97 precomputed checkpoints,
distributed as shown in Table~\ref{tab:selection-grid-size}. No candidate was
trained or modified during selection.

\begin{table}[h]
  \centering
  \small
  \begin{tabular}{lrrrr}
    \toprule
    Method & Rome & Baseball & AI & Total \\
    \midrule
    EMBER & 12 & 9 & 9 & 30 \\
    RMU & 8 & 8 & 8 & 24 \\
    SNMF & 4 & 3 & 3 & 10 \\
    RMU+EMBER & 8 & 8 & 8 & 24 \\
    SNMF+EMBER & 3 & 3 & 3 & 9 \\
    \midrule
    Total & 35 & 31 & 31 & 97 \\
    \bottomrule
  \end{tabular}
  \caption{Number of candidate checkpoints evaluated for each method and
  concept.}
  \label{tab:selection-grid-size}
\end{table}

Selection used 50 target questions, 50 neighboring questions, and the
same 50 SciQ questions for every concept. For each question, we scored all
four options with the character-normalized likelihood ranking defined in
Section~\ref{sec:evaluation}. We then calculated the normalized accuracies,
efficacy, preservation, and $H_{\mathrm{selection}}$ values defined in
Section~\ref{sec:evaluation}. The full model exceeded the four-choice
chance level in every selection group. Because the score is a nested harmonic
mean, a low target-efficacy or preservation component limits the final value.

We ranked candidates first by $H_{\mathrm{selection}}$, then by preservation,
and then by efficacy. Remaining exact ties were resolved by a fixed
method-specific order: smaller edit intensity for EMBER; smaller layer, band,
and $\alpha$ tuple for RMU; and a declared selection-mode and edited-side order
for SNMF. The tie-breaking rule was fixed without consulting the twins or any
test result. Table~\ref{tab:selected-checkpoints} records the resulting
configurations.

For each ensemble, the independently selected EMBER winner was fixed first,
after which the RMU or SNMF grid was applied to that edited checkpoint. SNMF
features were re-derived from the EMBER-edited model. Ensemble selection used
the same score and tie-breaking order as the corresponding standalone method.

\subsection{EMBER Semantic Feature Selection}
\label{app:ember-feature-selection}

EMBER used \texttt{google/gemma-4-12B-it} to judge whether candidate factors
represented the target concept, with a minimum confidence of 0.85. For Rome
and Baseball, the feature selected by the judge was also the unique factor
with the highest target-to-neutral ratio. The final edit grids therefore used
a numeric threshold that isolated exactly that previously judged feature,
avoiding repeated judge calls without changing the selected set. For AI,
some factors rejected by the judge had higher ratios than accepted factors,
so no single threshold could reproduce the semantic decision. The AI runs
therefore invoked the judge directly. The semantic selection criterion was
the same across concepts; only the mechanism used to reproduce its output in
the final grid differed.

\subsection{Selected Checkpoints}

\begin{center}
  \centering
  \scriptsize
  \resizebox{\columnwidth}{!}{%
  \begin{tabular}{lll}
    \toprule
    Concept & Method & Selected configuration \\
    \midrule
    Rome & EMBER & $\delta=200$ \\
         & RMU   & layer 6, high steering, $\alpha=10$ \\
         & SNMF  & ratio selection, output side$^{\dagger}$ \\
         & RMU+EMBER & layer 5, mid steering, $\alpha=100$$^{\ddagger}$ \\
         & SNMF+EMBER & ratio selection, input side \\
    \midrule
    Baseball & EMBER & $\delta=10$ \\
             & RMU   & layer 6, high steering, $\alpha=10$ \\
             & SNMF  & ratio selection, both sides$^{\dagger}$ \\
             & RMU+EMBER & layer 5, mid steering, $\alpha=100$$^{\ddagger}$ \\
             & SNMF+EMBER & ratio selection, both sides \\
    \midrule
    AI & EMBER & $\delta=500$ \\
       & RMU   & layer 6, high steering, $\alpha=10$ \\
       & SNMF  & ratio selection, input side \\
       & RMU+EMBER & layer 5, mid steering, $\alpha=10$$^{\ddagger}$ \\
       & SNMF+EMBER & ratio selection, input side \\
    \bottomrule
  \end{tabular}%
  }
  \captionof{table}{Configurations selected using $H_{\mathrm{selection}}$.
  Here, $\delta$ is EMBER's edit strength and $\alpha$ is RMU's retain weight.
  $^{\dagger}$ All SNMF candidates for this concept had zero efficacy and hence
  $H_{\mathrm{selection}}=0$; the displayed configuration was determined by
  the fixed tie-breaking order rather than a measured efficacy difference.
  $^{\ddagger}$ The selected RMU+EMBER checkpoint failed the engagement check on
  RMU's control direction, so it measures the selection rule rather than RMU's
  contribution; see Section~\ref{sec:results}.}
  \label{tab:selected-checkpoints}
\end{center}

The selection score and candidate ranking used only the selection questions
and did not use the concept-excluded twins. Held-out measurements and twin
comparisons did not affect the selected configurations. We computed
$H_{\mathrm{test}}$ on the test splits with the same formula and also computed
$H_{\mathrm{test}}$ for the twin, using the full model as the normalization
baseline, to provide a concept-exclusion reference for the erasure models.

The selection scorer used deterministic teacher-forced log probabilities in
FP32; no sampling or generation was involved.
